\subsection{A map of the existing approaches}
The methods surveyed here differ both in how they represent a structure and in how they turn that representation into a prediction. Keeping those decisions separate is useful when reading the literature. An elastic-network model gives its operator a physical interpretation under modeling assumptions. A supervised descriptor pipeline can use a related matrix simply as a source of input features. Similar-looking matrices therefore need not imply identical scientific claims.

\begin{table}[htbp]\centering\small
\caption{Conceptual map of the approaches discussed in the cited literature. This is a guide to their roles, not a new comparison of predictive accuracy.}
\begin{tabularx}{\linewidth}{>{\raggedright\arraybackslash}p{.23\linewidth}>{\raggedright\arraybackslash}X>{\raggedright\arraybackslash}X}\toprule
Approach & Structural representation & Quantity emphasized\\\midrule
Elastic networks & A residue contact network, with scalar or direction-dependent interactions & Collective fluctuation modes under the network model \citep{haliloglu1997,atilgan2001}.\\
Flexibility--rigidity indices & Distance-dependent sums around atoms or residues & Local interaction density across selected kernels and scales \citep{nguyen2016}.\\
Graph organization & Contacts and small graph configurations & The arrangement of contacts beyond a single count \citep{praznikar2023}.\\
Persistent homology & A nested family of complexes & The appearance and disappearance of topological classes \citep{xia2014}.\\
Persistent Laplacians & Operators associated with pairs of nested complexes & Kernel information together with positive spectral information \citep{wang2020,memoli2022,davies2023}.\\
Cellular sheaf spectra & Cellwise vector spaces and compatible comparison maps & Cohomology and weighted disagreement of local assignments \citep{hansen2019,wei2025}.\\\bottomrule
\end{tabularx}
\end{table}

Three successive questions organize this progression. Which points should interact? Which higher-dimensional cells should be included? How should assignments on different cells be compared? A contact graph answers the first question. A simplicial complex also answers the second. A cellular sheaf supplies the third ingredient. Persistence then asks how a compatible construction behaves across two members of a nested family. These choices are independent enough that they should be stated explicitly rather than hidden in a descriptor's name.

The distinction between geometry and topology will recur throughout the manuscript. Geometry includes distances, angles, and weights. Topology concerns connectivity and higher-dimensional holes, as represented by the chosen complex. Changing distances can alter a positive eigenvalue even when the complex and its Betti numbers stay unchanged. Conversely, adding a triangle can remove a one-dimensional topological class without changing any vertex or edge. A spectral representation can contain both kinds of information, but it should not call every spectral change a topological change.

The controlled comparison in this work places these representations inside one fixed learning setup. It asks whether the chosen summaries add useful predictive information beyond a specified geometric baseline. It does not compare optimized implementations of every method in the table. This narrower question makes it possible to connect the empirical outcome to the detailed operator analysis that follows.
